% !TEX program = lualatex
% resolucion.tex: resolución de la Secretaría (p. ej., representantes en
% comisiones tras un pleno). Ejemplo con nombres inventados.
\documentclass[11pt,a4paper]{article}
\usepackage{acta-consejo}
\usepackage{tabularx}
\usepackage{array}
\usepackage{colortbl}
% config.tex: datos fijos de los documentos (órgano, firmante y marca BORRADOR).
% Es lo único que hay que cambiar al empezar un curso nuevo.
\setorgano{el}{Consejo de Estudiantes}
\renewcommand{\firmante}{Nombre Apellidos}
\renewcommand{\cargofirmante}{Secretario \delorgano{} de la Facultad de Ciencias}
% Mientras el documento no sea definitivo, deja \borradortrue: añade la marca
% «BORRADOR» de fondo. Para la versión que se firma, \borradorfalse.
\borradortrue


\newcommand{\cab}[1]{\cellcolor{azulUZ}\textcolor{white}{\bfseries #1}}
\newcommand{\desierto}{\emph{[Desierto]} & \\\hline}

% \organo{Título}{filas}: tabla de un órgano con sus representantes.
% Cada fila: Nombre & Curso de incorporación\\\hline
\newcommand{\organoT}[2]{%
  \par\Needspace{5\baselineskip}\vspace{4pt}\noindent
  {\setlength{\parskip}{0pt}\setlength{\tabcolsep}{4pt}\renewcommand{\arraystretch}{1.05}%
  \begin{tabularx}{\linewidth}{|X|p{0.3\linewidth}|}
  \hline
  \multicolumn{2}{|l|}{\cab{#1}}\\\hline
  \cab{Representante} & \cab{Curso de incorporación}\\\hline
  #2
  \end{tabularx}}\par}

\begin{document}

\begin{center}
{\LARGE\bfseries\color{azulUZ} RESOLUCIÓN 01.2099: Representantes en comisiones}
\end{center}

En la presente resolución, la Secretaría ratifica a los representantes
designados por el Pleno \delorgano{} de la Facultad de Ciencias celebrado el
1 de enero de 2099.

\punto{Comisiones de la Facultad.}

\organoT{Comisión de Movilidad}{%
  Lucía Ejemplo Pérez & 2098/2099\\\hline
  \desierto}

\punto{Observaciones.}
\begin{itemize}
\item Queda desierta una plaza en la Comisión de Movilidad.
\end{itemize}

Para que conste a los efectos oportunos, se firma a fecha de la firma digital.
\firma

\end{document}
